\documentclass[11pt]{article}
\usepackage{ap-lit}
\title{Unit 1: Prose Fiction}
\author{Paul Wang}
\date{November 2026}
\docpart{Unit 1 \textperiodcentered\ Study Guide \textperiodcentered\ Short fiction 42--49\% of MCQ}

\begin{document}
\maketitle

\begin{abstract}
Short fiction is the largest share of the multiple choice, and a prose passage is the subject of the second essay. This guide takes the course's skill categories in the order you will use them on a passage: character (section 1), setting (2), plot and structure (3), the narrator (4), word choice, imagery, and symbols (5), and comparison and irony (6). Section 7 turns the toolkit into a plan for Question 2, the Prose Fiction Analysis essay. Examples are written for this guide unless a source is named. Each section ends with check questions; the key is on the last page. The Unit 1 Question Set applies all of this to three public-domain passages.
\end{abstract}

\scoringline{Check questions: \blank\ / \qtotal \qquad Target: 13+ \qquad Skill I miss most: \blank[6em]}

\section{Character}
\label{sec:1}

\subsection{How character is built}
\textbf{Direct characterization} states a trait (``She was a woman of mean understanding''). \textbf{Indirect characterization} shows it through speech, action, thought, appearance, and how others react, and it is what most questions are about. Ask of every character detail: \emph{what does this reveal that the narrator does not say?}

\begin{terms}
Round / flat & A round character has several, sometimes conflicting traits; a flat character is built around one. \\
Dynamic / static & A dynamic character changes in an important way; a static one does not. Change is often a change in \emph{understanding}, not behavior. \\
Foil & A character whose contrasting traits sharpen another's (a cautious sister beside a reckless one). \\
Motivation & What a character wants and why; conflicts arise when motives clash, within one person or between people. \\
Complex character & One who holds contradictory values, desires, or self-images at once. Essay prompts almost always ask about this. \\
\end{terms}

\begin{keyidea}[Read contrasts]
A character's values are clearest where two things collide: what she says versus what she does, how she sees herself versus how others see her, what she wants versus what she can admit wanting.
\end{keyidea}

\begin{q}{2}{1.1 Indirect characterization}{An action reveals the trait without stating it: indirect characterization.}
``He tipped the waiter generously, then spent the walk home complaining about the price of the meal.'' The sentence characterizes the man mainly through
\ch{direct statement by the narrator}{a contrast between his public action and his private complaint}{a description of his appearance}{the reactions of other characters}
\end{q}

\begin{q}{4}{1.1 Dynamic character}{A change in understanding makes a character dynamic.}
At the end of a story, a girl realizes that the aunt she has always admired is unhappy, though nothing in her life outwardly changes. The girl is best described as
\chh{flat}{a foil}{static}{dynamic}
\end{q}

\section{Setting}
\label{sec:2}

Setting is time, place, and social world: the season, the room, the class and customs that press on a character. On the exam, setting is rarely just background. It can \textbf{mirror} a character's inner state, \textbf{contrast} with it (a wedding in a downpour), \textbf{confine} a character (narrow houses, locked rooms), or \textbf{reveal} values (what a family keeps on display). \textbf{Atmosphere} or \textbf{mood} is the feeling a setting creates in the reader.

\begin{trap}
Do not assume the weather always matches the character's mood. When the setting is bright and the character is miserable, the \emph{contrast} is usually the point: the world does not share her feeling, or she is out of step with it.
\end{trap}

\begin{q}{3}{1.2 Setting}{The cramped room externalizes how confined her life feels: setting mirrors character.}
A story describes a woman's ``narrow room, where the ceiling sloped so low she could stand upright only at the window.'' The setting most likely serves to
\ch{establish the historical period}{introduce a mystery to be solved later}{suggest the limits that press on her life, and the window as her one outlet}{show that she is wealthy}
\end{q}

\begin{q}{1}{1.2 Setting as contrast}{A festive setting set against private grief isolates the character.}
A widow walks through a crowded street fair on the day of her husband's funeral. The setting most likely emphasizes
\ch{her isolation from the life around her}{the happiness she shares with the crowd}{the economic history of the town}{the reliability of the narrator}
\end{q}

\section{Plot and structure}
\label{sec:3}

\subsection{Arrangement}
\textbf{Plot} is the arrangement of events, not just the events themselves. Writers choose where to begin (\textbf{in medias res}, in the middle of action), whether to go back (\textbf{flashback}) or hint ahead (\textbf{foreshadowing}), and how fast to move: \textbf{scene} slows time to a moment-by-moment account; \textbf{summary} covers months in a sentence. A change in pacing marks what matters.

\subsection{Shifts and endings}
Most passages turn: from description to reflection, from past to present, from a character's view to the narrator's, from confidence to doubt. An \textbf{epiphany} is a sudden moment of insight, often near a story's end (James Joyce made it a signature). Many modern stories end \textbf{open}, without resolving the conflict, so the reader must weigh what the character's realization means.

\begin{terms}
Exposition & Background on characters and situation. \\
Conflict & Internal (within a character) or external (with another person, society, nature). \\
Climax / turning point & The moment of greatest tension, where the conflict tips. \\
Juxtaposition & Two scenes, images, or ideas placed side by side to invite comparison. \\
Frame & A story told within another story or by a narrator looking back. \\
\end{terms}

\begin{q}{2}{1.3 Pacing}{Slowing into detailed scene marks a moment's importance.}
A story covers ten years of a marriage in one paragraph, then spends three pages on a single dinner. The change in pacing mainly signals that
\ch{the narrator is unreliable}{the dinner is a decisive moment}{the earlier years were happier}{the story is told out of order}
\end{q}

\begin{q}{4}{1.3 Epiphany}{A sudden realization that reframes experience = epiphany.}
At a story's end, a boy standing in an emptying hall suddenly sees his devotion as vanity. This moment is best described as
\chh{exposition}{a flashback}{foreshadowing}{an epiphany}
\end{q}

\section{The narrator}
\label{sec:4}

The narrator is the most tested element on the exam. The narrator is a voice inside the text, never simply the author.

\begin{compare}{Point of view}{What the narrator can do}{What to ask}
First person & Tells as ``I,'' limited to what that person knows and admits. & What does this narrator not see, or not say? \\
Retrospective first person & An older self narrates the younger self's experience. Two perspectives at once: the \emph{experiencing self} and the \emph{narrating self}. & Where does the older voice judge, ironize, or mourn the younger? \\
Third person limited & ``He'' or ``she,'' close to one character's perceptions. & How close is the narration to that character's mind? \\
Third person omniscient & Moves among minds and can comment on them. & What does the narrator's commentary add, and with what attitude? \\
Objective & Reports only what could be seen and heard. & What do gestures and silences reveal? \\
\end{compare}

\subsection{Distance and voice}
\textbf{Narrative distance} is how close the telling stands to a character's consciousness. \textbf{Free indirect discourse} (or style) blends a character's thoughts into third-person narration without quotation marks: ``She would not cry. Not here, not in front of all of them.'' The words are the character's, in the narrator's grammar. When the narrator steps outside a character to comment or generalize, the narration is \textbf{intrusive}. A narrator is \textbf{unreliable} when we have reason to doubt the account, because of youth, bias, self-deception, or limited knowledge.

\begin{trap}
In free indirect discourse, an exclamation or a judgment in the narration often belongs to the character, not the narrator. ``What a splendid day it was!'' may tell you how the character feels, and the rest of the passage may show that the day is not splendid at all.
\end{trap}

\begin{q}{1}{1.4 Retrospective narrator}{An adult looking back can judge the boy's feelings while still describing them.}
In a story narrated by a man recalling his boyhood, he calls his early infatuation ``my foolish blood.'' The word ``foolish'' most clearly reveals
\ch{the older narrator's judgment of his younger self}{the boy's awareness at the time that his love was foolish}{the girl's opinion of the boy}{an omniscient narrator's view of all the characters}
\end{q}

\begin{q}{3}{1.4 Free indirect discourse}{Character's thought in third-person grammar, no quotation marks.}
``Mrs.\ Hale looked at the clock. Late again. Well, let him eat it cold, then.'' The last two sentences are best described as
\ch{direct dialogue}{the objective narrator's report}{free indirect discourse, rendering Mrs.\ Hale's thoughts}{an authorial generalization}
\end{q}

\begin{q}{2}{1.4 Intrusive narrator}{A generalizing comment outside any character's mind = intrusive narration.}
A narrator pauses a scene to remark, ``How few of us, at nineteen, know what we are refusing.'' This is an example of
\chh{dialogue}{intrusive narration}{an epiphany}{an objective narrator}
\end{q}

\section{Word choice, imagery, and symbols}
\label{sec:5}

\textbf{Diction} is word choice. Two words can denote the same thing and \textbf{connote} very different things: a house can be ``modest,'' ``plain,'' or ``shabby.'' \textbf{Detail} is what the writer chooses to include. \textbf{Imagery} is language that appeals to the senses. A \textbf{motif} is a recurring image or idea; a \textbf{symbol} is an object, image, or action that carries meaning beyond itself, usually signaled by repetition or emphasis.

\textbf{Tone} is the narrator's or speaker's attitude toward the subject; \textbf{mood} is the feeling created in the reader. \textbf{Syntax} (sentence structure) carries tone too: a run of short sentences can sound clipped or stunned; a long \textbf{cumulative} sentence piles on detail; a \textbf{periodic} sentence delays its main clause for suspense; repetition can sound obsessive, ritual, or insistent.

\begin{terms}
Tone words & Wistful, sardonic, reverent, detached, wry, elegiac, indignant, ambivalent, tender, mock-heroic. \\
Ambiguity & A word, image, or ending that supports more than one reading, often on purpose. \\
Understatement & Saying less than the situation warrants, often for irony or restraint. \\
Hyperbole & Deliberate exaggeration. \\
\end{terms}

\begin{q}{4}{1.5 Connotation}{``Lurked'' connotes threat; the others are neutral or positive.}
Which verb gives the most threatening connotation to ``The old house \blank[2.5em] at the end of the street''?
\chh{stood}{rested}{waited}{lurked}
\end{q}

\begin{q}{1}{1.5 Syntax}{A sequence of short sentences after long ones sounds abrupt and stunned.}
After two long, flowing sentences describing a picnic, a story says: ``Then the telegram. She read it twice. She folded it.'' The change in syntax mainly creates
\ch{a sense of abrupt shock and numbness}{a relaxed, leisurely mood}{comic exaggeration}{a flashback to the past}
\end{q}

\begin{q}{3}{1.5 Symbol}{Repeated, emphasized object linked to a value = symbol.}
A worn fur collar appears at the beginning and end of a story; the main character talks to it as a friend and, after being mocked, hears it ``crying.'' The collar most likely functions as
\ch{a detail with no larger meaning}{evidence of the character's wealth}{a symbol of the character's fragile self-image}{a flashback device}
\end{q}

\section{Comparison and irony}
\label{sec:6}

\begin{terms}
Simile & Comparison with ``like'' or ``as.'' \\
Metaphor & Implied comparison: one thing spoken of as another. An \textbf{extended metaphor} develops one comparison over several sentences. \\
Personification & Human qualities given to things or ideas (``the houses gazed at one another''). \\
Allusion & Reference to a well-known text, person, or event (a ``chalice'' recalling the Grail quest). \\
Verbal irony & Saying the opposite of, or something different from, what is meant. \\
Situational irony & An outcome opposite to what was expected. \\
Dramatic irony & The reader knows something a character does not. \\
\end{terms}

The question to ask of any comparison is \emph{what qualities transfer}. If a boy's body is ``like a harp,'' then the girl's words are what play it: he is an instrument, moved from outside. Comparisons tell you how a character or narrator understands the world.

\begin{q}{2}{1.6 Comparison}{Personified houses with ``faces'' that ``gazed'' make the street seem to judge.}
``The houses, conscious of decent lives within them, gazed at one another with brown imperturbable faces.'' The personification chiefly suggests a neighborhood that is
\ch{lively and welcoming}{smugly respectable and unmoved}{dangerous and violent}{under construction}
\end{q}

\begin{q}{4}{1.6 Irony}{The reader knows what the character does not: dramatic irony.}
A reader knows from an earlier scene that the letter a character is waiting for was burned, while the character continues to watch the post each morning. This is
\chh{an allusion}{verbal irony}{situational irony}{dramatic irony}
\end{q}

\section{Question 2: the Prose Fiction Analysis essay}
\label{sec:7}

The prompt gives a passage of roughly 500--700 words, a sentence of context, and a task in this form: \emph{analyze how the author uses literary elements and techniques to convey} a complex character, relationship, or experience. The word to underline is the one after ``convey''; your thesis must say something about it.

\begin{enumerate}[leftmargin=1.6em]
\item \textbf{Read the prompt}, then the passage, once through for the situation and narrator.
\item \textbf{Find the complexity}: the two forces in tension (hope and dread; pride and shame; what the character believes and what the narrator shows).
\item \textbf{Map the movement}: where does the passage begin, where does it shift, and where does it end? A chronological plan that follows that movement is often the clearest line of reasoning.
\item \textbf{Choose evidence} from at least three parts of the passage and at least two kinds of technique (narration, imagery, syntax, setting, dialogue).
\item \textbf{Write the thesis} with the complexity in it, then argue it paragraph by paragraph.
\end{enumerate}

\begin{keyidea}[Thesis frame for Q2]
Through [technique] and [technique], [author] portrays [character] as [first quality] yet [second quality], revealing [what the tension shows].
\end{keyidea}

\begin{q}{1}{1.7 Q2 thesis}{The prompt's object (a complex relationship) must be interpreted in the thesis.}
A prompt asks how the author conveys ``the complex relationship between a daughter and her father.'' Which thesis best responds?
\ch{Through the daughter's practiced silences and her one outburst, the author shows a love that has learned to disguise itself as obedience until the moment it cannot.}{The author uses dialogue, imagery, and setting to describe the relationship.}{The daughter and her father have a complex relationship.}{The passage is about a daughter and her father at dinner.}
\end{q}

\begin{q}{3}{1.7 Line of reasoning}{Following the passage's movement and its shift gives a clear line of reasoning.}
A passage moves from a character's eager anticipation to her quiet humiliation. Which plan best organizes a Q2 essay on it?
\ch{One paragraph on every simile, then one on every metaphor}{A summary of the plot followed by a list of devices}{One paragraph on how the opening builds her anticipation, one on the turn, and one on how the ending registers her humiliation}{A paragraph on the author's biography}
\end{q}

\makekey
\end{document}
